\documentclass[../DoAn.tex]{subfiles}
\begin{document}
\section{Tổng kết}
Mặc dù gặp khó khăn với phần nghiệp vụ hệ thống, nhưng nhờ góp ý của thầy Nguyễn Ngọc Duy, tụi em đã triển khai được phần nào chức năng của hệ thống, đảm bảo cơ bản về mặt chức năng nên có.
\subsection{Kết quả đạt được}
\begin{itemize}
	\item Phân tích và thiết kế cơ sở dữ liệu
	\item Nắm được các ngôn ngữ lập trình
	\item Khảo sát và nắm được sơ lược về quy trình của dự án \textit{Hệ thống phát hiện người mất tích dựa trên dữ liệu ảnh chụp}
	\item Triển khai được cơ sở dữ liệu trên dịch vụ AWS của Amazon
	\item Tìm hiểu thêm kĩ thuật trong mô hình học sâu CNN.
	\item Dự án bao gồm các chức năng cơ bản như: quản lý, thống kê, bảo mật theo quyền của tài khoản, \dots
	\item Đạt độ chính xác 96\% với dataset Pin Face Recognition, và 94\% với dataset LFW.
\end{itemize}
\subsection{Hạn chế}
Tốc độ phản hồi của dự án còn chậm. Còn một số điểm trong phần nghiệp vụ vẫn chưa sát với thực tế như tích hợp được CCTV trong thực tế để thực hiện truy vấn thêm dữ liệu. Mô hình học sâu
\section{Hướng phát triển}
Đây là đề tài mang tính thực tiển cao, có thể áp dụng vào thực tế với bối cảnh các vụ mất tích hiện nay xảy ra càng nhiều. Ứng dụng thêm công nghệ phát triển hiện nay như AI có thể hỗ trợ được phần nào gánh nặng trong công tác tìm kiếm, các nghiệp vụ triển khai nhanh chóng.

Hướng tiếp cận tương lai cho dự án có thể là:
\begin{itemize}
	\item Cải thiện được tốc độ phản hồi của dự án.
	\item Triển khai ra được cộng đồng cũng như môi trường thực tế.
\end{itemize}
\end{document}
